\chapter*{Abstract}
\addcontentsline{toc}{chapter}{Abstract}
\begin{SingleSpace}
This thesis presents the design and implementation of Lectura, an AI-driven educational platform that transforms unstructured multimedia content—such as protracted lecture videos and dense PDF documents—into highly consolidated, interactive study materials. As digital educational resources grow exponentially, university students face an acute cognitive burden and information overload. To address this, Lectura leverages Google's Gemini 3 Flash Large Language Model to automatically generate structured Cornell notes, context-aware multiple-choice quizzes, algorithmic spaced repetition flashcards, and visually organized presentation slides natively exportable to PDF and PPTX.

The platform is engineered as a resilient, cloud-native full-stack web application. It features a React 18 frontend with inline WYSIWYG editing, seamlessly connected to a highly concurrent Go (Golang) backend. By utilizing a robust worker pool architecture, the system efficiently manages long-running asynchronous AI inference tasks while ensuring real-time responsiveness via WebSockets. Ultimately, by integrating automated content ingestion with established pedagogical frameworks like active recall and Bloom's Taxonomy, Lectura significantly reduces the administrative burden of manual note-taking, providing a unified and high-performance environment for autonomous learning.
\end{SingleSpace}

\chapter*{Аңдатпа}
\addcontentsline{toc}{chapter}{Аңдатпа}
\begin{SingleSpace}
Бұл дипломдық жұмыс ұзақ бейнелекциялар мен күрделі PDF құжаттары сияқты құрылымдалмаған мультимедиялық мазмұнды жоғары топтастырылған, интерактивті оқу материалдарына айналдыратын жасанды интеллектке негізделген Lectura білім беру платформасын әзірлеу мен енгізуді ұсынады. Цифрлық білім беру ресурстарының өсуімен студенттер үлкен когнитивті жүктеме мен ақпараттық шамадан тыс жүктемеге ұшырайды. Осы мәселені шешу үшін жүйе Google Gemini 3 Flash үлкен тілдік моделін (LLM) пайдаланып, құрылымдалған Корнелл конспектілерін, тесттерді, интервалдық қайталау флэш-карталарын және PDF пен PPTX форматтарына экспортталатын визуалды презентацияларды автоматты түрде жасайды.

Платформа сенімді, бұлтқа бағытталған толық стек веб-қосымшасы ретінде жасалған. Ол inline өңдеу мүмкіндігі бар React 18 фронтендінен және параллельді өңдеуді қамтамасыз ететін Go (Golang) бәкендінен тұрады. Жұмысшы пулы (worker pool) архитектурасын қолдану арқылы жүйе ұзақ уақытты қажет ететін асинхронды ЖИ тапсырмаларын тиімді басқарады. Автоматтандырылған мазмұнды өңдеуді дәлелденген педагогикалық әдістермен біріктіру арқылы Lectura оқу процесіндегі рутиндық жұмысты айтарлықтай азайтып, автономды оқыту үшін бірыңғай, өнімділігі жоғары ортаны қамтамасыз етеді.
\end{SingleSpace}

\chapter*{Аннотация}
\addcontentsline{toc}{chapter}{Аннотация}
\begin{SingleSpace}
В данной дипломной работе представлено проектирование и реализация Lectura — образовательной платформы на базе искусственного интеллекта, которая преобразует неструктурированный мультимедийный контент, такой как долгие видеолекции и объемные PDF-документы, в консолидированные интерактивные учебные материалы. В условиях экспоненциального роста цифровых ресурсов студенты сталкиваются с информационной перегрузкой. Для решения этой проблемы Lectura использует большую языковую модель Google Gemini 3 Flash, автоматически генерируя структурированные конспекты Корнелла, контекстно-зависимые тесты, флэш-карточки для интервального повторения и визуально оформленные презентации с возможностью экспорта в PDF и PPTX.

Платформа реализована как отказоустойчивое облачное веб-приложение. Она включает фронтенд на React 18 с инлайн-редактором (WYSIWYG) и высококонкурентный бэкенд на Go (Golang). Благодаря архитектуре пула воркеров (worker pool), система эффективно управляет длительными асинхронными задачами нейросети, сохраняя отзывчивость интерфейса. Объединяя автоматизированную обработку данных с такими педагогическими методиками, как интервальное повторение и активное вспоминание, Lectura значительно снижает рутинную нагрузку при подготовке к занятиям, предоставляя единую высокопроизводительную среду для автономного обучения.
\end{SingleSpace}
